\scanpage{80}
$D=D_0=D_{\delta_0}$ autour de $x$ dans le sens de
l'orientation donnée, jusqu'en $D_1$ passant
par $s_\alpha$ (cette opération est anodine,
on aurait pu commencer par un point
$\delta_1$ voisin de $a_\alpha$). Puis suivant le long de
$a_\alpha$ suivant l'orientation induite par $F$
--- ce qui signifie que $D_1$ pivote autour de
$s_\alpha$, dans un sens qu'il faudra expliciter
tantôt (déterminé par l'orientation de $D_1$ déduite
de celle de $D_0$ par trans. pal., et la rive
privilégiée dont on dispose. Notons que l'
orientation \inserted{$\omega_\alpha$} de $\Delta_\alpha$ qui porte $a_\alpha$, celle $\unsure{\omega_{\delta_1}}$ du point
$\delta_1$ sur $\Delta_\alpha$, et la rive $\rho_\alpha$ de $\Delta_\alpha$ déterminée
par $F$, forment un syst. compatible au sens de
Stokes, et qu'il s'ensuit que les
orientations correspondantes en l'élément de
contact $(s_\alpha,D_1)$, $\omega=\omega_\alpha^\vee\in\uline{\omega}(X,s_\alpha)$, $\omega_{\delta_1}^\vee=\omega_{D_1}$, et
$\rho_\alpha^\vee=$ rive de balayage, forment un système
\uline{anticompatible}, i.e. l'orientation $\omega_{D_1}$ de $D_1$
est \uline{opposée} à celle induite par $\omega$ sur la
rive de balayage (cf figure). Ceci peut \unsure{se|ainsi}
se propager le long de $a_\alpha$ \inserted{dans la direction $\omega_\alpha$} ce qui fait tourner
$D_1$ autour de $s_\alpha$ dans la \unsure{direction},
indiquée, jusqu'au moment où
on heurte un premier sommet
\inserted{(qui sera $s_\beta$)} --- qui est aussi celui où le
point mobile $\delta$ sur $a_\alpha$ arrive
au sommet $\delta_2$, rencontre de
\deleted{$\unsure{\Delta_\alpha\text{ et}}$} $a_\alpha$ et $a_\beta$ i.e. de $\Delta_\alpha$ et $\Delta_\beta$
[Bien sûr, \inserted{entre les arcs de} \deleted{\unsure{en cours du}}
\deleted{déplacement de $D_1$, \unsure{arêtes}}
spécialisations correspondantes $a'_\alpha$ et
$a'_\beta$, il peut y avoir, dans un
sens comme dans l'autre
($\omega_D$ ou l'orientation opposée) bien
d'autres arcs de spécialisation, mais
correspondants à des sommets $s_\gamma$
qui ne sont pas atteints au cours
du balayage]. A partir de $s_\beta,D_\beta$
on continue comme précédemment.
\margin{79}
